\section{Conclusion}\label{sec:conclusion}

This paper isolates a narrow Bayesian pricing question: how much does integrating uncertainty in historically estimated volatility change a WTI Average Price Option value relative to plugging in the posterior mean? The answer is governed by posterior dispersion and pricing curvature. Synthetic experiments confirm this mechanism and show that posterior integration can become economically relevant when information is weak and the pricing map is sufficiently curved.

The observed WTI sample lies in a different region. Posterior-integrated and posterior-mean point values are very close, yet the posterior distribution of theoretical prices remains materially wider. The distinction matters: a small \(C_{PI}-C_{PM}\) is a statement about nonlinear averaging of the posterior, not about the absence of parameter-induced price uncertainty.

The larger empirical margin is the volatility information supplied to the pricing map. Across the tested forward benchmarks, prior option-implied volatility states predict APO settlements much better than the historical-volatility inputs. That finding survives recent-window, EWMA, and horizon-matched GARCH comparisons; it does not require flexible smile fitting; and it remains visible when the volatility source is moved to an independent vanilla-WTI option market. These results support an option-informed versus tested-historical comparison, not a pure structural decomposition between physical and risk-neutral volatility.

The contribution is therefore diagnostic rather than a new option-pricing theory. Posterior integration matters when dispersion meets curvature, but in this sample improving the volatility information is much more consequential for point pricing than more exact integration of the historical-volatility posterior. The evidence is limited by the small number of distinct market dates, uneven liquidity, and the maintained one-factor pricing model. Richer commodity dynamics can be studied from this benchmark without conflating model enrichment with the narrower effect of Bayesian posterior integration.
